% GENERATED FILE. Edit canonical lessons, then run npm run generate:books.
% canonical-source-hash: fnv1a64:c32bd67294ee11d7
% canonical-lessons: BN-C464-gombhir, BN-C464-sadasidhe, BN-C464-khutkhute, BN-C464-udasin, BN-C464-maniye-neoya

\chapter{Grave and Easygoing}
\label{ch:bn-grave-and-easygoing}
\hlchaptermodality{464}

\begin{chapteropening}
\textbf{By the end of this chapter:} \emph{I can say whether someone is serious, simple, fussy or indifferent.}

Complete the last lesson of chapter 464: I can say whether someone is serious, simple, fussy or indifferent.
\end{chapteropening}

\section[\texorpdfstring{\bn{গম্ভীর} (gômbhīr)}{গম্ভীর (gômbhīr)}]{\bn{গম্ভীর} (gômbhīr) — grave; serious}

\label{lesson:BN-C464-gombhir}



\begin{quote}
Before the new one: say the Bengali for to tease, then the Bengali for to tickle.
\end{quote}

\subsection*{You'll want to know: \bn{গম্ভীর}}
\textbf{\bn{গম্ভীর}} — \emph{gômbhīr} — \textquotedblleft{}grave; serious\textquotedblright{}.

\begin{quote}
\bn{খবরটা} \bn{শুনে} \bn{বাবা} \bn{গম্ভীর} \bn{হয়ে} \bn{গেলেন।}
\end{quote}

\begin{quote}
\emph{khôbôrṭā shune bābā gômbhīr hôye gelen.}
\end{quote}

\begin{quote}
Father went grave when he heard the news.
\end{quote}

\begin{cousinweb}
\textbf{\bn{গম্ভীর}} (\emph{gômbhīr}) is a tatsama, Sanskrit \emph{gambhīra}, \textquotedblleft{}deep\textquotedblright{} (CDIAL 4031), which Monier-Williams gives as a form of \emph{gabhīra}, the source of \bn{গভীর} (taught), deep; Bengali uses it for a deep, serious manner.
\end{cousinweb}

\begin{grammarlens}[title={What you've built}]
One more word for describing what you lived and felt.
\end{grammarlens}

\subsection*{Your turn}
\begin{itemize}
\raggedright
  \item \emph{Say it:} \emph{gômbhīr}
  \item \emph{Say it:} \emph{gômbhīr}, once more
  \item \emph{Say it:} \emph{khôbôrṭā shune bābā gômbhīr hôye gelen.} — Father went grave when he heard the news.
  \item \emph{Recall:} say the Bengali for to tease, then the Bengali for to tickle, then say \emph{gômbhīr} again
\end{itemize}

\begin{tcolorbox}[breakable,colback=teal!4,colframe=teal!35!black,title={Before you move on}]
What does \bn{গম্ভীর} mean? (\textquotedblleft{}Grave; serious\textquotedblright{}.) Say it once more.
\end{tcolorbox}
\section[\texorpdfstring{\bn{সাদাসিধে} (sādāsidhe)}{সাদাসিধে (sādāsidhe)}]{\bn{সাদাসিধে} (sādāsidhe) — plain and simple; unpretentious}

\label{lesson:BN-C464-sadasidhe}



\begin{quote}
Before the new one: say the Bengali for to tickle, then the Bengali for grave.
\end{quote}

\subsection*{You'll want to know: \bn{সাদাসিধে}}
\textbf{\bn{সাদাসিধে}} — \emph{sādāsidhe} — \textquotedblleft{}plain and simple; unpretentious\textquotedblright{}.

\begin{quote}
\bn{দাদুর} \bn{অনেক} \bn{টাকা}, \bn{কিন্তু} \bn{তিনি} \bn{সাদাসিধে} \bn{জীবন} \bn{কাটান।}
\end{quote}

\begin{quote}
\emph{dādur ônek ṭākā, kintu tini sādāsidhe jībôn kāṭān.}
\end{quote}

\begin{quote}
Grandfather has plenty of money, but he lives a plain and simple life.
\end{quote}

\begin{cousinweb}
\textbf{\bn{সাদাসিধে}} (\emph{sādāsidhe}) pairs \bn{সাদা} (taught), white and so plain, with \emph{sidhe}, the Bengali \emph{sidhā}, \textquotedblleft{}straight, honest\textquotedblright{}, which CDIAL lists under Sanskrit \emph{siddha}, \textquotedblleft{}accomplished\textquotedblright{} (CDIAL 13401).
\end{cousinweb}

\begin{grammarlens}[title={What you've built}]
One more word for describing what you lived and felt.
\end{grammarlens}

\subsection*{Your turn}
\begin{itemize}
\raggedright
  \item \emph{Say it:} \emph{sādāsidhe}
  \item \emph{Say it:} \emph{sādāsidhe}, once more
  \item \emph{Say it:} \emph{dādur ônek ṭākā, kintu tini sādāsidhe jībôn kāṭān.} — Grandfather has plenty of money, but he lives a plain and simple life.
  \item \emph{Recall:} say the Bengali for to tickle, then the Bengali for grave, then say \emph{sādāsidhe} again
\end{itemize}

\begin{tcolorbox}[breakable,colback=teal!4,colframe=teal!35!black,title={Before you move on}]
What does \bn{সাদাসিধে} mean? (\textquotedblleft{}Plain and simple; unpretentious\textquotedblright{}.) Say it once more.
\end{tcolorbox}
\section[\texorpdfstring{\bn{খুঁতখুঁতে} (khũtkhũte)}{খুঁতখুঁতে (khũtkhũte)}]{\bn{খুঁতখুঁতে} (khũtkhũte) — fussy; hard to please}

\label{lesson:BN-C464-khutkhute}



\begin{quote}
Before the new one: say the Bengali for grave, then the Bengali for unpretentious.
\end{quote}

\subsection*{You'll want to know: \bn{খুঁতখুঁতে}}
\textbf{\bn{খুঁতখুঁতে}} — \emph{khũtkhũte} — \textquotedblleft{}fussy; hard to please\textquotedblright{}.

\begin{quote}
\bn{খাবার} \bn{নিয়ে} \bn{অরুণ} \bn{খুঁতখুঁতে।}
\end{quote}

\begin{quote}
\emph{khābār niye ôrun khũtkhũte.}
\end{quote}

\begin{quote}
Arun is fussy about food.
\end{quote}

\begin{cousinweb}
\textbf{\bn{খুঁতখুঁতে}} (\emph{khũtkhũte}) doubles \emph{khũt}, \textquotedblleft{}a flaw\textquotedblright{}: someone who finds a flaw in everything. The origin of \emph{khũt} is not settled.
\end{cousinweb}

\begin{grammarlens}[title={What you've built}]
One more word for describing what you lived and felt.
\end{grammarlens}

\subsection*{Your turn}
\begin{itemize}
\raggedright
  \item \emph{Say it:} \emph{khũtkhũte}
  \item \emph{Say it:} \emph{khũtkhũte}, once more
  \item \emph{Say it:} \emph{khābār niye ôrun khũtkhũte.} — Arun is fussy about food.
  \item \emph{Recall:} say the Bengali for grave, then the Bengali for unpretentious, then say \emph{khũtkhũte} again
\end{itemize}

\begin{tcolorbox}[breakable,colback=teal!4,colframe=teal!35!black,title={Before you move on}]
What does \bn{খুঁতখুঁতে} mean? (\textquotedblleft{}Fussy; hard to please\textquotedblright{}.) Say it once more.
\end{tcolorbox}
\section[\texorpdfstring{\bn{উদাসীন} (udāsīn)}{উদাসীন (udāsīn)}]{\bn{উদাসীন} (udāsīn) — indifferent; careless (of something)}

\label{lesson:BN-C464-udasin}



\begin{quote}
Before the new one: say the Bengali for unpretentious, then the Bengali for fussy.
\end{quote}

\subsection*{You'll want to know: \bn{উদাসীন}}
\textbf{\bn{উদাসীন}} — \emph{udāsīn} — \textquotedblleft{}indifferent; careless (of something)\textquotedblright{}.

\begin{quote}
\bn{স্বাস্থ্য} \bn{নিয়ে} \bn{অরুণ} \bn{উদাসীন।}
\end{quote}

\begin{quote}
\emph{svāsthyô niye ôrun udāsīn.}
\end{quote}

\begin{quote}
Arun does not care about his health.
\end{quote}

\begin{cousinweb}
\textbf{\bn{উদাসীন}} (\emph{udāsīn}) is a tatsama, Sanskrit \emph{udāsīna}, \textquotedblleft{}sitting apart, indifferent\textquotedblright{} (Monier-Williams).
\end{cousinweb}

\begin{grammarlens}[title={What you've built}]
One more word for describing what you lived and felt.
\end{grammarlens}

\subsection*{Your turn}
\begin{itemize}
\raggedright
  \item \emph{Say it:} \emph{udāsīn}
  \item \emph{Say it:} \emph{udāsīn}, once more
  \item \emph{Say it:} \emph{svāsthyô niye ôrun udāsīn.} — Arun does not care about his health.
  \item \emph{Recall:} say the Bengali for unpretentious, then the Bengali for fussy, then say \emph{udāsīn} again
\end{itemize}

\begin{tcolorbox}[breakable,colback=teal!4,colframe=teal!35!black,title={Before you move on}]
What does \bn{উদাসীন} mean? (\textquotedblleft{}Indifferent; careless (of something)\textquotedblright{}.) Say it once more.
\end{tcolorbox}
\section[\texorpdfstring{\bn{মানিয়ে} \bn{নেওয়া} (māniye neoyā)}{মানিয়ে নেওয়া (māniye neoyā)}]{\bn{মানিয়ে} \bn{নেওয়া} (māniye neoyā) — to adjust; to adapt (to people or a place)}

\label{lesson:BN-C464-maniye-neoya}



\begin{quote}
Before the new one: say the Bengali for fussy, then the Bengali for indifferent.
\end{quote}

\subsection*{You'll want to know: \bn{মানিয়ে} \bn{নেওয়া}}
\textbf{\bn{মানিয়ে} \bn{নেওয়া}} — \emph{māniye neoyā} — \textquotedblleft{}to adjust; to adapt (to people or a place)\textquotedblright{}.

\begin{quote}
\bn{নতুন} \bn{শহরে} \bn{মীরা} \bn{তাড়াতাড়ি} \bn{মানিয়ে} \bn{নিল।}
\end{quote}

\begin{quote}
\emph{nôtun shôhôre mīrā tāṛātāṛi māniye nilô.}
\end{quote}

\begin{quote}
Mira adapted quickly to the new city.
\end{quote}

Past (he/she …): \bn{তিনি} \bn{নতুন} \bn{চাকরিতে} \bn{মানিয়ে} \bn{নিলেন}, \emph{tini nôtun chākrite māniye nilen}.

\begin{cousinweb}
\textbf{\bn{মানিয়ে} \bn{নেওয়া}} (\emph{māniye neoyā}) is \emph{māniye}, the -\bn{এ} form of \emph{mānāno}, \textquotedblleft{}to fit, to suit\textquotedblright{}, from the verb \emph{mānā} that CDIAL lists for Bengali under Sanskrit \emph{manyate} (CDIAL 9857), with \emph{neoyā}, \textquotedblleft{}to take\textquotedblright{} (CDIAL 6966).
\end{cousinweb}

\begin{grammarlens}[title={What you've built}]
One more word for describing what you lived and felt.
\end{grammarlens}

\subsection*{Your turn}
\begin{itemize}
\raggedright
  \item \emph{Say it:} \emph{māniye neoyā}
  \item \emph{Say it:} \emph{māniye neoyā}, once more
  \item \emph{Say it:} \emph{nôtun shôhôre mīrā tāṛātāṛi māniye nilô.} — Mira adapted quickly to the new city.
  \item \emph{Recall:} say the Bengali for fussy, then the Bengali for indifferent, then say \emph{māniye neoyā} again
\end{itemize}

\begin{tcolorbox}[breakable,colback=teal!4,colframe=teal!35!black,title={Before you move on}]
What does \bn{মানিয়ে} \bn{নেওয়া} mean? (\textquotedblleft{}To adjust; to adapt (to people or a place)\textquotedblright{}.) Say it once more.
\end{tcolorbox}
